\documentclass[tikz,border=4pt]{standalone}
\usepackage{ctex}
\usepackage{amsmath}
\usepackage{pgfplots}
\pgfplotsset{compat=1.18}

% p10-08a: 铅垂高法图——抛物线y=-x^2+2x+3，A(-1,0), B(3,0), C(0,3)，P为动点
% 铅垂高 h = |y_P - y_Q|，底BC，Q是BC上横坐标与P相同的点
\begin{document}
\begin{tikzpicture}
\begin{axis}[
  width=9cm, height=8cm,
  axis lines=center,
  xlabel={$x$}, ylabel={$y$},
  every axis x label/.style={at={(axis cs:4.0,0)}, anchor=west},
  every axis y label/.style={at={(axis cs:0,4.5)}, anchor=south},
  xmin=-2.0, xmax=4.0,
  ymin=-1.0, ymax=4.5,
  xtick={-2,-1,1,2,3},
  ytick={1,2,3,4},
  tick style={color=black},
  ticklabel style={font=\small},
  clip=false,
]
  % 抛物线 y = -x^2 + 2x + 3
  \addplot[blue, thick, domain=-1.3:3.3, samples=100] {-x^2 + 2*x + 3};
  \node[right, blue] at (axis cs:2.8, 1.8) {$y=-x^2+2x+3$};

  % 点 A(-1,0), B(3,0), C(0,3)
  \addplot[only marks, mark=*, mark size=3pt] coordinates {(-1,0) (3,0) (0,3)};
  \node[below] at (axis cs:-1,0) {$A(-1,0)$};
  \node[below] at (axis cs:3,0) {$B(3,0)$};
  \node[left] at (axis cs:0,3) {$C(0,3)$};

  % 底边 BC（红色）
  \draw[red, thick] (axis cs:0,3) -- (axis cs:3,0);
  \node[above right, red] at (axis cs:1.5,1.5) {\small 底$BC$};

  % 动点 P(3/2, 15/4)（面积最大时的点）
  \addplot[only marks, mark=*, mark size=3pt, red] coordinates {(1.5, 3.75)};
  \node[above] at (axis cs:1.5, 3.75) {$P$};

  % 底直线BC方程 y=-x+3，在x=1.5处Q(1.5, 1.5)
  \addplot[only marks, mark=*, mark size=2.5pt, gray] coordinates {(1.5, 1.5)};
  \node[below right, gray] at (axis cs:1.5, 1.5) {$Q$};

  % 铅垂高（竖线P到Q）
  \draw[red, thick, <->] (axis cs:1.5, 1.5) -- (axis cs:1.5, 3.75) node[right, midway, red] {铅垂高$h$};

  % 水平宽（BC两端横坐标之差）
  \draw[gray, dashed] (axis cs:0,0) -- (axis cs:0,3);
  \draw[gray, dashed] (axis cs:3,0) -- (axis cs:3,0);
  \draw[gray, <->, dashed] (axis cs:0,-0.6) -- (axis cs:3,-0.6) node[midway, below] {\small 水平宽$|x_B-x_C|=3$};

  % x=1.5辅助虚线
  \draw[gray, dashed] (axis cs:1.5, 0) -- (axis cs:1.5, 1.5);
\end{axis}
\end{tikzpicture}
\end{document}
